\chapter{Ordering Ramen in Japanese}

% ================= INTRO =================
\begin{convobox}{1}
    \noindent
    \jpkana{みなさん}{Minasan} \jpkana{こんにちは、}{konnichiwa,} 
    \jpkana{たなかです。}{Tanaka desu.}

    \vspace{1.5em}
    \noindent{\Large \color{articolor} \textbf{Arti:} Halo semuanya, saya Tanaka.}
\end{convobox}

\begin{convobox}{2}
    \noindent
    \jpword{きょう}{今日}{Kyou} \jpkana{は、}{wa,} 
    \jpkana{ラーメン}{raamen}\jpword{や}{屋}{ya} \jpkana{で}{de} 
    \jpword{ちゅうもん}{注文}{chuumon} \jpkana{するときの}{suru toki no} 
    \jpword{にほんご}{日本語}{Nihongo} \jpkana{を}{o} 
    \jpword{いっしょ}{一緒}{issho} \jpkana{に}{ni} 
    \jpword{べんきょう}{勉強}{benkyou} \jpkana{しましょう。}{shimashou.}

    \vspace{1.5em}
    \noindent{\Large \color{articolor} \textbf{Arti:} Hari ini, mari kita belajar bersama bahasa Jepang saat memesan di kedai ramen.}
\end{convobox}

\begin{convobox}{3}
    \noindent
    \jpword{とつぜん}{突然}{Totsuzen} \jpkana{ですが、}{desu ga,} 
    \jpkana{みなさんは}{minasan wa} 
    \jpkana{ラーメン}{raamen} \jpkana{が}{ga} 
    \jpword{す}{好}{su}\jpkana{きですか？}{ki desu ka?} 
    \jpword{わたし}{私}{Watashi} \jpkana{は}{wa} 
    \jpword{だいす}{大好}{daisu}\jpkana{きです。}{ki desu.}

    \vspace{1.5em}
    \noindent{\Large \color{articolor} \textbf{Arti:} Tiba-tiba (Oiya), apakah kalian suka ramen? Kalau saya sangat suka.}
\end{convobox}

\begin{convobox}{4}
    \noindent
    \jpword{にほん}{日本}{Nihon} \jpkana{に}{ni} 
    \jpword{き}{来}{ki}\jpkana{たら、}{tara,} 
    \jpkana{ラーメン}{raamen}\jpword{や}{屋}{ya}\jpword{めぐ}{巡}{megu}\jpkana{りを}{ri o} 
    \jpkana{したいという}{shitai to iu} 
    \jpword{かた}{方}{kata} \jpkana{も}{mo} 
    \jpword{おお}{多}{oo}\jpkana{いと}{i to} 
    \jpword{おも}{思}{omo}\jpkana{います。}{imasu.}

    \vspace{1.5em}
    \noindent{\Large \color{articolor} \textbf{Arti:} Jika datang ke Jepang, aku rasa banyak orang yang ingin berwisata kuliner (berkeliling) kedai ramen.}
\end{convobox}

\begin{convobox}{5}
    \noindent
    \jpkana{でも、}{Demo,} 
    \jpkana{「}{"}\jpword{にほんご}{日本語}{Nihongo} \jpkana{で}{de} 
    \jpkana{ラーメン}{raamen} \jpkana{を}{o} 
    \jpword{ちゅうもん}{注文}{chuumon} \jpkana{するのは、}{suru no wa,} 
    \jpword{すこ}{少}{suko}\jpkana{し}{shi} 
    \jpword{むずか}{難}{muzuka}\jpkana{しそうだなぁ」}{shisou da naa"}

    \vspace{1.5em}
    \noindent{\Large \color{articolor} \textbf{Arti:} Tapi, "Memesan ramen dalam bahasa Jepang, sepertinya sedikit sulit ya..."}
\end{convobox}

\begin{convobox}{6}
    \noindent
    \jpkana{...と}{...to} 
    \jpword{おも}{思}{omo}\jpkana{っている}{tte iru} 
    \jpword{かた}{方}{kata} \jpkana{も}{mo} 
    \jpkana{いるのではないでしょうか。}{iru no de wa nai deshou ka.}

    \vspace{1.5em}
    \noindent{\Large \color{articolor} \textbf{Arti:} ...mungkin ada juga yang berpikir seperti itu.}
\end{convobox}

\begin{convobox}{7}
    \noindent
    \jpkana{でも、}{Demo,} 
    \jpword{あんしん}{安心}{anshin} \jpkana{してください。}{shite kudasai.} 
    \jpkana{この}{Kono} \jpword{どうが}{動画}{douga} \jpkana{を}{o} 
    \jpword{み}{観}{mi}\jpkana{れば、}{reba,} 
    \jpkana{ラーメン}{raamen} \jpkana{の}{no} 
    \jpword{ちゅうもん}{注文}{chuumon} \jpword{ほうほう}{方法}{houhou} \jpkana{が}{ga} 
    \jpword{わ}{分}{wa}\jpkana{かります。}{karimasu.}

    \vspace{1.5em}
    \noindent{\Large \color{articolor} \textbf{Arti:} Tapi, tenang saja (jangan khawatir). Kalau menonton video ini, kamu akan mengerti cara memesan ramen.}
\end{convobox}

\begin{convobox}{8}
    \noindent
    \jpkana{さらに、}{Sara ni,} 
    \jpkana{トッピングを}{toppingu o} 
    \jpword{ついか}{追加}{tsuika} \jpkana{したり、}{shitari,} 
    \jpword{めん}{麺}{men} \jpkana{を}{o} 
    \jpword{この}{この}{kono}\jpkana{みの}{mi no} 
    \jpkana{かたさにしてもらうように}{katasa ni shite morau you ni} 
    \jpword{つた}{伝}{tsuta}\jpkana{えたりして、}{etari shite,}

    \vspace{1.5em}
    \noindent{\Large \color{articolor} \textbf{Arti:} Terlebih lagi, (kamu bisa) menambah topping, menyampaikan tingkat kekerasan mi sesuai seleramu,}
\end{convobox}

\begin{convobox}{9}
    \noindent
    \jpkana{もっと}{motto} \jpkana{ラーメン}{raamen} \jpkana{を}{o} 
    \jpkana{おいしく}{oishiku} 
    \jpword{た}{食}{ta}\jpkana{べることも}{beru koto mo} 
    \jpkana{できるようになりますよ。}{dekiru you ni narimasu yo.}

    \vspace{1.5em}
    \noindent{\Large \color{articolor} \textbf{Arti:} sehingga kamu akan bisa memakan ramen dengan lebih lezat lagi lho.}
\end{convobox}

\begin{convobox}{10}
    \noindent
    \jpkana{それでは}{Sore dewa} 
    \jpword{さっそく}{早速}{sassoku} 
    \jpword{い}{行}{i}\jpkana{ってみましょう！}{tte mimashou!}

    \vspace{1.5em}
    \noindent{\Large \color{articolor} \textbf{Arti:} Kalau begitu, ayo langsung kita mulai!}
\end{convobox}

% ================= TITLE: CONVERSATION =================
\vspace{2em}
\begin{tcolorbox}[colback=boxframe, colframe=boxframe, rounded corners, boxrule=0pt, arc=3mm, left=2mm, top=2mm, bottom=2mm]
    \Large \bfseries \color{white} 🍜 1. Conversation (Percakapan Dasar Memesan Ramen)
\end{tcolorbox}
\vspace{1em}

\begin{convobox}{11}
    \noindent
    \jpword{まず}{先ず}{Mazu} \jpword{いちど}{一度}{ichido} 
    \jpkana{ラーメン}{raamen}\jpword{や}{屋}{ya} \jpkana{での}{de no} 
    \jpword{かいわ}{会話}{kaiwa} \jpkana{を}{o} 
    \jpword{しぜん}{自然}{shizen} \jpkana{な}{na} 
    \jpkana{スピードで}{supiido de} 
    \jpword{き}{聞}{ki}\jpkana{いてみましょう。}{ite mimashou.}

    \vspace{1.5em}
    \noindent{\Large \color{articolor} \textbf{Arti:} Pertama-tama, mari kita dengarkan sekali percakapan di kedai ramen dengan kecepatan natural.}
\end{convobox}

\begin{convobox}{12}
    \noindent
    \jpkana{いらっしゃいませ、}{Irasshaimase,} 
    \jpword{なんめい}{何名}{nan-mei}\jpword{さま}{様}{sama} \jpkana{ですか？}{desu ka?}

    \vspace{1.5em}
    \noindent{\Large \color{articolor} \textbf{Arti:} Selamat datang, untuk berapa orang?}
\end{convobox}

\begin{convobox}{13}
    \noindent
    \jpword{ひとり}{1人}{Hitori} \jpkana{です。}{desu.}

    \vspace{1.5em}
    \noindent{\Large \color{articolor} \textbf{Arti:} 1 orang.}
\end{convobox}

\begin{convobox}{14}
    \noindent
    \jpkana{こちらへ}{Kochira e} \jpkana{どうぞ。}{douzo.}

    \vspace{1.5em}
    \noindent{\Large \color{articolor} \textbf{Arti:} Silakan ke sebelah sini.}
\end{convobox}

\begin{convobox}{15}
    \noindent
    \jpkana{しょうゆラーメン}{Shouyu raamen} 
    \jpkana{お}{o}\jpword{ねが}{願}{nega}\jpkana{いします。}{ishimasu.}

    \vspace{1.5em}
    \noindent{\Large \color{articolor} \textbf{Arti:} Tolong pesan ramen rasa Shoyu (kecap asin).}
\end{convobox}

\begin{convobox}{16}
    \noindent
    \jpword{めん}{麺}{Men} \jpkana{の}{no} 
    \jpkana{かたさは}{katasa wa} 
    \jpkana{どうなさいますか？}{dou nasaimasu ka?}

    \vspace{1.5em}
    \noindent{\Large \color{articolor} \textbf{Arti:} Tingkat kematangan (kekerasan) mi-nya bagaimana?}
\end{convobox}

\begin{convobox}{17}
    \noindent
    \jpword{めん}{麺}{Men} 
    \jpkana{かためで}{katame de} 
    \jpkana{お}{o}\jpword{ねが}{願}{nega}\jpkana{いします。}{ishimasu.}

    \vspace{1.5em}
    \noindent{\Large \color{articolor} \textbf{Arti:} Tolong mi-nya agak keras (setengah matang).}
\end{convobox}

\begin{convobox}{18}
    \noindent
    \jpkana{かしこまりました。}{Kashikomarimashita.}

    \vspace{1.5em}
    \noindent{\Large \color{articolor} \textbf{Arti:} Baik, saya mengerti.}
\end{convobox}

\begin{convobox}{19}
    \noindent
    \jpkana{トッピングは}{Toppingu wa} 
    \jpkana{いかがですか？}{ikaga desu ka?}

    \vspace{1.5em}
    \noindent{\Large \color{articolor} \textbf{Arti:} Bagaimana dengan tambahan topping-nya?}
\end{convobox}

\begin{convobox}{20}
    \noindent
    \jpword{あじたま}{味玉}{Ajitama} 
    \jpkana{トッピング}{toppingu} 
    \jpkana{お}{o}\jpword{ねが}{願}{nega}\jpkana{いします。}{ishimasu.}

    \vspace{1.5em}
    \noindent{\Large \color{articolor} \textbf{Arti:} Tolong tambah topping telur rebus (Ajitama).}
\end{convobox}

\begin{convobox}{21}
    \noindent
    \jpkana{かしこまりました。}{Kashikomarimashita.}

    \vspace{1.5em}
    \noindent{\Large \color{articolor} \textbf{Arti:} Baik.}
\end{convobox}

\begin{convobox}{22}
    \noindent
    \jpkana{それでは、}{Sore dewa,} 
    \jpword{じゅんばん}{順番}{junban} \jpkana{に}{ni} 
    \jpword{せつめい}{説明}{setsumei} \jpkana{しますね。}{shimasu ne.}

    \vspace{1.5em}
    \noindent{\Large \color{articolor} \textbf{Arti:} Kalau begitu, akan saya jelaskan secara berurutan ya.}
\end{convobox}

\begin{lessonbox}{Tips Belanja untuk Anak \& Orang Tua: Jurus "Kore, Onegaishimasu"}
    \textbf{TIPS UMUM:} Saat masuk ke restoran mana saja di Jepang, pelayan akan langsung menyambut dengan \textit{Irasshaimase} dan membawa menu. 
    \\
    \textbf{JURUS RAHASIA:} Kalau kamu belum tahu cara membaca menunya, jangan khawatir! Cukup \textbf{tunjuk gambarnya} di buku menu dan ucapkan: \textbf{"Kore, onegaishimasu"} (Tolong pesan yang INI). Begitu mudah!
\end{lessonbox}

% ================= TITLE: PHRASE 1 =================
\vspace{2em}
\begin{tcolorbox}[colback=boxframe, colframe=boxframe, rounded corners, boxrule=0pt, arc=3mm, left=2mm, top=2mm, bottom=2mm]
    \Large \bfseries \color{white} 👥 2. Phrase 1 (Cara Menghitung Jumlah Orang)
\end{tcolorbox}
\vspace{1em}

\begin{convobox}{23}
    \noindent
    \jpkana{「いらっしゃいませ、}{"Irasshaimase,} 
    \jpword{なんめい}{何名}{nan-mei}\jpword{さま}{様}{sama} \jpkana{ですか？」}{desu ka?"}

    \vspace{1.5em}
    \noindent{\Large \color{articolor} \textbf{Arti:} "Selamat datang, untuk berapa orang?"}
\end{convobox}

\begin{convobox}{24}
    \noindent
    \jpkana{「}{"}\jpword{ひとり}{1人}{Hitori} \jpkana{です。」}{desu."}

    \vspace{1.5em}
    \noindent{\Large \color{articolor} \textbf{Arti:} "1 orang."}
\end{convobox}

\begin{convobox}{25}
    \noindent
    \jpkana{お}{O}\jpword{みせ}{店}{mise} \jpkana{の}{no} 
    \jpword{ひと}{人}{hito} \jpkana{に}{ni} 
    \jpword{にんずう}{人数}{ninzuu} \jpkana{を}{o} 
    \jpword{つた}{伝}{tsuta}\jpkana{える}{eru} 
    \jpkana{フレーズです。}{fureezu desu.}

    \vspace{1.5em}
    \noindent{\Large \color{articolor} \textbf{Arti:} Ini adalah frasa untuk menyampaikan jumlah orang kepada staf toko/restoran.}
\end{convobox}

\begin{convobox}{26}
    \noindent
    \jpword{にんずう}{人数}{Ninzuu} \jpkana{は、}{wa,} 
    \jpkana{このように}{kono you ni} 
    \jpword{い}{言}{i}\jpkana{います。}{imasu.}

    \vspace{1.5em}
    \noindent{\Large \color{articolor} \textbf{Arti:} Jumlah orang, diucapkan seperti ini.}
\end{convobox}

\begin{convobox}{27}
    \noindent
    \jpword{ひとり}{1人}{Hitori} / \jpword{ふたり}{2人}{Futari} / \jpword{さんにん}{3人}{San-nin} / \jpword{よにん}{4人}{Yo-nin} / \jpword{ごにん}{5人}{Go-nin}

    \vspace{1.5em}
    \noindent{\Large \color{articolor} \textbf{Arti:} 1 orang / 2 orang / 3 orang / 4 orang / 5 orang}
\end{convobox}

\begin{convobox}{28}
    \noindent
    \jpkana{このように}{Kono you ni} 
    \jpword{ゆび}{指}{yubi} \jpkana{も}{mo} 
    \jpword{いっしょ}{一緒}{issho} \jpkana{に}{ni} 
    \jpword{だ}{出}{da}\jpkana{すと、}{su to,} 
    \jpkana{お}{o}\jpword{みせ}{店}{mise} \jpkana{の}{no} 
    \jpword{ひと}{人}{hito} \jpkana{に}{ni} 
    \jpword{わ}{分}{wa}\jpkana{かりやすいです。}{kariyasui desu.}

    \vspace{1.5em}
    \noindent{\Large \color{articolor} \textbf{Arti:} Jika kamu juga menunjukkan jari (sesuai jumlah) seperti ini, staf toko akan lebih mudah memahaminya.}
\end{convobox}

\begin{lessonbox}{Hafalan Wajib: Menghitung Anggota Keluarga!}
    \textbf{Orang Tua \& Anak:} Di Jepang, menghitung orang itu sedikit unik. 1 dan 2 orang panggilannya spesial!
    \begin{itemize}
        \item Sendiri = \textbf{Hitori}
        \item Berdua = \textbf{Futari}
        \item Bertiga = \textbf{San-nin} (Tinggal tambah "nin" di belakang angka)
        \item Berempat = \textbf{Yo-nin} (Bukan Yon-nin lho ya!)
    \end{itemize}
    \textbf{Aksi:} Saat pelayan tanya \textit{"Nan-mei sama desu ka?"}, angkat jarimu tinggi-tinggi dan ucapkan jumlah keluargamu dengan bangga! "Yo-nin desu!" (4 orang).
\end{lessonbox}

\begin{convobox}{29}
    \noindent
    \jpkana{「しょうゆラーメン}{"Shouyu raamen} 
    \jpkana{お}{o}\jpword{ねが}{願}{nega}\jpkana{いします。」}{ishimasu."}

    \vspace{1.5em}
    \noindent{\Large \color{articolor} \textbf{Arti:} "Tolong ramen rasa shoyu."}
\end{convobox}

\begin{convobox}{30}
    \noindent
    \jpkana{ラーメンの}{Raamen no} 
    \jpword{しゅるい}{種類}{shurui} \jpkana{を}{o} 
    \jpword{つた}{伝}{tsuta}\jpkana{える}{eru} 
    \jpkana{フレーズです。}{fureezu desu.}

    \vspace{1.5em}
    \noindent{\Large \color{articolor} \textbf{Arti:} Ini adalah frasa untuk menyebutkan jenis ramen.}
\end{convobox}

\begin{convobox}{31}
    \noindent
    \jpkana{ラーメンの}{Raamen no} \jpword{しゅるい}{種類}{shurui} \jpkana{には、}{ni wa,} 
    \jpword{たと}{例}{tato}\jpkana{えば}{eba} 
    \jpkana{このようなものが}{kono you na mono ga} 
    \jpkana{あります。}{arimasu.}

    \vspace{1.5em}
    \noindent{\Large \color{articolor} \textbf{Arti:} Untuk jenis ramen, misalnya ada yang seperti ini.}
\end{convobox}

\begin{convobox}{32}
    \noindent
    \jpkana{しょうゆラーメン}{Shouyu raamen} / 
    \jpkana{とんこつラーメン}{Tonkotsu raamen} / 
    \jpword{しお}{塩}{Shio} \jpkana{ラーメン}{raamen} / 
    \jpkana{みそラーメン}{Miso raamen}

    \vspace{1.5em}
    \noindent{\Large \color{articolor} \textbf{Arti:} Ramen Shoyu (Kecap asin) / Ramen Tonkotsu (Kaldu tulang) / Ramen Shio (Garam) / Ramen Miso (Pasta kedelai)}
\end{convobox}

\begin{convobox}{33}
    \noindent
    \jpkana{ひとつだけでなく、}{Hitotsu dake de naku,} 
    \jpkana{いくつか}{ikutsu ka} 
    \jpword{ちゅうもん}{注文}{chuumon} \jpkana{するときは、}{suru toki wa,} 
    \jpkana{このように}{kono you ni} 
    \jpword{い}{言}{i}\jpkana{うことができます。}{u koto ga dekimasu.}

    \vspace{1.5em}
    \noindent{\Large \color{articolor} \textbf{Arti:} Jika memesan bukan cuma satu (tapi beberapa), kamu bisa mengucapkannya seperti ini.}
\end{convobox}

\begin{convobox}{34}
    \noindent
    \jpkana{「しょうゆラーメン}{"Shouyu raamen} \jpkana{ひとつ}{hitotsu} \jpkana{と、}{to,} 
    \jpkana{とんこつラーメン}{tonkotsu raamen} \jpkana{ふたつ}{futatsu} 
    \jpkana{お}{o}\jpword{ねが}{願}{nega}\jpkana{いします。」}{ishimasu."}

    \vspace{1.5em}
    \noindent{\Large \color{articolor} \textbf{Arti:} "Tolong ramen shoyu 1 (mangkuk), dan ramen tonkotsu 2 (mangkuk)."}
\end{convobox}

\begin{convobox}{35}
    \noindent
    \jpkana{「しょうゆラーメン}{"Shouyu raamen} \jpkana{と、}{to,} 
    \jpword{しお}{塩}{shio}\jpkana{ラーメンと、}{raamen to,} 
    \jpkana{みそラーメン}{miso raamen} 
    \jpkana{お}{o}\jpword{ねが}{願}{nega}\jpkana{いします。」}{ishimasu."}

    \vspace{1.5em}
    \noindent{\Large \color{articolor} \textbf{Arti:} "Tolong ramen shoyu, dan ramen shio, dan ramen miso."}
\end{convobox}

\begin{lessonbox}{Hafalan Wajib: Menghitung Benda (Porsi/Mangkuk)!}
    Kalau tadi menghitung orang, sekarang menghitung benda bulat atau porsi mangkuk!
    \begin{itemize}
        \item 1 Porsi = \textbf{Hitotsu}
        \item 2 Porsi = \textbf{Futatsu}
        \item 3 Porsi = \textbf{Mittsu}
    \end{itemize}
    \textbf{Rumusnya:} [Nama Makanan] + [Jumlah] + [Onegaishimasu]. \\
    Contoh: "Ramen Hitotsu, Onegaishimasu!" (Tolong Ramen 1 porsi!)
\end{lessonbox}

% ================= TITLE: PHRASE 3 =================
\vspace{2em}
\begin{tcolorbox}[colback=boxframe, colframe=boxframe, rounded corners, boxrule=0pt, arc=3mm, left=2mm, top=2mm, bottom=2mm]
    \Large \bfseries \color{white} 🥢 4. Phrase 3 (Kekerasan Mi)
\end{tcolorbox}
\vspace{1em}

\begin{convobox}{36}
    \noindent
    \jpkana{「}{"}\jpword{めん}{麺}{Men} \jpkana{かためで}{katame de} 
    \jpkana{お}{o}\jpword{ねが}{願}{nega}\jpkana{いします。」}{ishimasu."}

    \vspace{1.5em}
    \noindent{\Large \color{articolor} \textbf{Arti:} "Tolong mi-nya agak keras (setengah matang)."}
\end{convobox}

\begin{convobox}{37}
    \noindent
    \jpkana{お}{O}\jpword{みせ}{店}{mise} \jpkana{によっては、}{ni yotte wa,} 
    \jpkana{このように}{kono you ni} 
    \jpword{めん}{麺}{men} \jpkana{の}{no} 
    \jpkana{かたさを}{katasa o} 
    \jpword{えら}{選}{era}\jpkana{ぶことができます。}{bu koto ga dekimasu.}

    \vspace{1.5em}
    \noindent{\Large \color{articolor} \textbf{Arti:} Tergantung tokonya, kamu bisa memilih tingkat kekerasan/kematangan mi seperti ini.}
\end{convobox}

\begin{convobox}{38}
    \noindent
    \jpkana{かためが}{Katame ga} \jpkana{いい}{ii} \jpword{とき}{時}{toki}\jpkana{は、}{wa,} 
    \jpkana{「}{"}\jpword{めん}{麺}{men} \jpkana{かためで}{katame de} 
    \jpkana{お}{o}\jpword{ねが}{願}{nega}\jpkana{いします。」}{ishimasu."}

    \vspace{1.5em}
    \noindent{\Large \color{articolor} \textbf{Arti:} Jika ingin yang agak keras (al dente), "Tolong mi-nya katame."}
\end{convobox}

\begin{convobox}{39}
    \noindent
    \jpkana{やわらかめが}{Yawarakame ga} \jpkana{いい}{ii} \jpword{とき}{時}{toki}\jpkana{は、}{wa,} 
    \jpkana{「}{"}\jpword{めん}{麺}{men} \jpkana{やわらかめで}{yawarakame de} 
    \jpkana{お}{o}\jpword{ねが}{願}{nega}\jpkana{いします。」}{ishimasu."}

    \vspace{1.5em}
    \noindent{\Large \color{articolor} \textbf{Arti:} Jika ingin yang lembek/lunak, "Tolong mi-nya yawarakame."}
\end{convobox}

\begin{convobox}{40}
    \noindent
    \jpword{ふつう}{普通}{Futsuu} \jpkana{の}{no} 
    \jpkana{かたさが}{katasa ga} \jpkana{いい}{ii} \jpword{とき}{時}{toki}\jpkana{は、}{wa,} 
    \jpkana{「}{"}\jpword{ふつう}{普通}{futsuu} \jpkana{で}{de} 
    \jpkana{お}{o}\jpword{ねが}{願}{nega}\jpkana{いします。」}{ishimasu."}

    \vspace{1.5em}
    \noindent{\Large \color{articolor} \textbf{Arti:} Jika ingin yang normal/standar, "Tolong yang futsuu (normal)."}
\end{convobox}

\begin{convobox}{41}
    \noindent
    \jpkana{ちなみに、}{Chinami ni,} 
    \jpword{わたし}{私}{watashi} \jpkana{は}{wa} 
    \jpword{めん}{麺}{men} \jpkana{かためが}{katame ga} 
    \jpword{す}{好}{su}\jpkana{きです。}{ki desu.}

    \vspace{1.5em}
    \noindent{\Large \color{articolor} \textbf{Arti:} Ngomong-ngomong, aku suka mi yang katame (agak keras/kenyal).}
\end{convobox}

\begin{lessonbox}{Aksi Jagoan Ramen: Pilih Kekerasan Mi!}
    Orang Jepang sangat serius soal mi ramen! Kalau pelayan menanyakan tentang *Katasa* (Kekerasan), kamu harus siap menjawab:
    \begin{itemize}
        \item \textbf{Katame:} Agak keras/kenyal
        \item \textbf{Futsuu:} Biasa/Standar.
        \item \textbf{Yawarakame:} Lembut/Lunak (Cocok untuk anak-anak kecil!).
    \end{itemize}
\end{lessonbox}

% ================= TITLE: PHRASE 4 =================
\vspace{2em}
\begin{tcolorbox}[colback=boxframe, colframe=boxframe, rounded corners, boxrule=0pt, arc=3mm, left=2mm, top=2mm, bottom=2mm]
    \Large \bfseries \color{white} 🥚 5. Phrase 4 (Tambah Topping)
\end{tcolorbox}
\vspace{1em}

\begin{convobox}{42}
    \noindent
    \jpkana{「}{"}\jpword{あじたま}{味玉}{Ajitama} \jpkana{トッピング}{toppingu} 
    \jpkana{お}{o}\jpword{ねが}{願}{nega}\jpkana{いします。」}{ishimasu."}

    \vspace{1.5em}
    \noindent{\Large \color{articolor} \textbf{Arti:} "Tolong topping telur rebus (kecap)."}
\end{convobox}

\begin{convobox}{43}
    \noindent
    \jpkana{ラーメンには、}{Raamen ni wa,} 
    \jpword{じぶん}{自分}{jibun} \jpkana{の}{no} 
    \jpword{この}{この}{kono}\jpkana{みに}{mi ni} 
    \jpword{あ}{合}{a}\jpkana{わせて}{wasete} 
    \jpkana{いろいろな}{iroiro na} 
    \jpkana{トッピングが}{toppingu ga} \jpkana{できます。}{dekimasu.}

    \vspace{1.5em}
    \noindent{\Large \color{articolor} \textbf{Arti:} Di ramen, kamu bisa melakukan bermacam topping sesuai dengan selera dirimu sendiri.}
\end{convobox}

\begin{convobox}{44}
    \noindent
    \jpkana{チャーシュー}{Chaashuu} (Daging iris) / \jpword{あじたま}{味玉}{Ajitama} (Telur rebus) / \jpkana{メンマ}{Menma} (Rebung) / \jpkana{ねぎ}{Negi} (Daun bawang) / \jpkana{のり}{Nori} (Rumput laut) / \jpkana{にんにく}{Ninniku} (Bawang putih) / \jpword{ほうれんそう}{ほうれん草}{Hourensou} (Bayam)
\end{convobox}

\begin{convobox}{45}
    \noindent
    \jpword{ちゅうもん}{注文}{Chuumon} \jpkana{するときは、}{suru toki wa,} 
    \jpkana{このように}{kono you ni} 
    \jpword{い}{言}{i}\jpkana{います。}{imasu.}

    \vspace{1.5em}
    \noindent{\Large \color{articolor} \textbf{Arti:} Saat memesan, ucapkan seperti ini.}
\end{convobox}

\begin{convobox}{46}
    \noindent
    \jpkana{「チャーシュー}{"Chaashuu} \jpkana{トッピング}{toppingu} 
    \jpkana{お}{o}\jpword{ねが}{願}{nega}\jpkana{いします。」}{ishimasu."}

    \vspace{1.5em}
    \noindent{\Large \color{articolor} \textbf{Arti:} "Tolong topping chaashuu."}
\end{convobox}

\begin{convobox}{47}
    \noindent
    \jpkana{「}{"}\jpword{あじたま}{味玉}{Ajitama} \jpkana{と、}{to,} 
    \jpkana{メンマと、}{menma to,} 
    \jpkana{にんにく}{ninniku} 
    \jpkana{トッピング}{toppingu} 
    \jpkana{お}{o}\jpword{ねが}{願}{nega}\jpkana{いします。」}{ishimasu."}

    \vspace{1.5em}
    \noindent{\Large \color{articolor} \textbf{Arti:} "Tolong topping ajitama, menma, dan ninniku."}
\end{convobox}

\begin{lessonbox}{Aksi Jagoan Ramen: Partikel "To" (Dan)}
    Untuk menggabungkan banyak topping, gunakan partikel \textbf{と (to)} yang artinya "Dan".\\
    Contoh: Ajitama \textbf{to}, Nori \textbf{to}, Negi \textbf{to}... onegaishimasu!
\end{lessonbox}

% ================= TITLE: TIPS =================
\vspace{2em}
\begin{tcolorbox}[colback=boxframe, colframe=boxframe, rounded corners, boxrule=0pt, arc=3mm, left=2mm, top=2mm, bottom=2mm]
    \Large \bfseries \color{white} 💡 6. Tips Rahasia Membayar: Mesin Tiket (Kenbaiki)
\end{tcolorbox}
\vspace{1em}

\begin{convobox}{48}
    \noindent
    \jpword{こんかい}{今回}{Konkai} \jpkana{は、}{wa,} 
    \jpword{てんいん}{店員}{tenin}\jpkana{さんに}{san ni} 
    \jpword{ちょくせつ}{直接}{chokusetsu} 
    \jpword{ちゅうもん}{注文}{chuumon} \jpkana{する}{suru} 
    \jpword{ほうほう}{方法}{houhou} \jpkana{を}{o} 
    \jpword{つた}{伝}{tsuta}\jpkana{えましたが、}{emashita ga,}

    \vspace{1.5em}
    \noindent{\Large \color{articolor} \textbf{Arti:} Kali ini, aku mengajarkan cara memesan langsung ke kasir/pelayan, namun}
\end{convobox}

\begin{convobox}{49}
    \noindent
    \jpkana{お}{O}\jpword{みせ}{店}{mise} \jpkana{によっては、}{ni yotte wa,} 
    \jpkana{このような}{kono you na} 
    \jpword{きかい}{機械}{kikai} \jpkana{で}{de} 
    \jpword{ちゅうもん}{注文}{chuumon} \jpkana{することもあります。}{suru koto mo arimasu.}

    \vspace{1.5em}
    \noindent{\Large \color{articolor} \textbf{Arti:} Tergantung tokonya, ada kalanya kamu memesan lewat mesin seperti ini.}
\end{convobox}

\begin{convobox}{50}
    \noindent
    \jpkana{この}{Kono} \jpword{きかい}{機械}{kikai} \jpkana{のことを、}{no koto o,} 
    \jpkana{「}{"}\jpword{けんばいき}{券売機}{kenbaiki}\jpkana{」と}{" to} 
    \jpword{い}{言}{i}\jpkana{います。}{imasu.}

    \vspace{1.5em}
    \noindent{\Large \color{articolor} \textbf{Arti:} Mesin ini, disebut "Kenbaiki" (Mesin penjual tiket otomatis).}
\end{convobox}

\begin{convobox}{51}
    \noindent
    \jpword{けんばいき}{券売機}{Kenbaiki} \jpkana{で}{de} 
    \jpword{ちゅうもん}{注文}{chuumon} \jpkana{するときは、}{suru toki wa,} 
    \jpword{まず}{先ず}{mazu} \jpkana{お}{o}\jpword{かね}{金}{kane} \jpkana{を}{o} 
    \jpword{い}{入}{i}\jpkana{れて、}{rete,}

    \vspace{1.5em}
    \noindent{\Large \color{articolor} \textbf{Arti:} Saat memesan di mesin, pertama masukkan uangnya,}
\end{convobox}

\begin{convobox}{52}
    \noindent
    \jpword{ちゅうもん}{注文}{chuumon} \jpkana{したい}{shitai} 
    \jpkana{ラーメンや}{raamen ya} 
    \jpkana{トッピングの}{toppingu no} 
    \jpkana{ボタンを}{botan o} \jpword{お}{押}{o}\jpkana{して、}{shite,}

    \vspace{1.5em}
    \noindent{\Large \color{articolor} \textbf{Arti:} lalu pencet tombol ramen atau topping yang ingin dipesan,}
\end{convobox}

\begin{convobox}{53}
    \noindent
    \jpword{で}{出}{de}\jpkana{てきた}{tekita} 
    \jpword{かみ}{紙}{kami} \jpkana{を}{o} 
    \jpkana{お}{o}\jpword{みせ}{店}{mise} \jpkana{の}{no} 
    \jpword{ひと}{人}{hito} \jpkana{に}{ni} 
    \jpword{わた}{渡}{wata}\jpkana{します。}{shimasu.}

    \vspace{1.5em}
    \noindent{\Large \color{articolor} \textbf{Arti:} lalu serahkan kertas (tiket) yang keluar itu kepada staf toko.}
\end{convobox}

\begin{convobox}{54}
    \noindent
    \jpkana{あとは、}{Ato wa,} 
    \jpkana{ラーメンが}{raamen ga} 
    \jpword{く}{来}{ku}\jpkana{るのを}{ru no o} 
    \jpword{ま}{待}{ma}\jpkana{つだけです。}{tsu dake desu.}

    \vspace{1.5em}
    \noindent{\Large \color{articolor} \textbf{Arti:} Setelahnya, kamu tinggal menunggu ramennya datang.}
\end{convobox}

\begin{lessonbox}{Aksi Penting Kenbaiki (Mesin Tiket) 🚨}
    \textbf{AWAS TERBALIK!} Ini kesalahan turis paling umum! \\
    Di Indonesia atau mesin minuman, kita tekan tombol dulu baru masukin uang kan? Di mesin ramen Jepang (\textbf{Kenbaiki}), kamu harus \textbf{MASUKKAN UANG DULU}, barulah lampu-lampu tombol menunya menyala untuk bisa ditekan! Setelah tiketnya keluar, berikan tiketnya ke pelayan tanpa perlu bicara apa-apa. Gampang banget!
\end{lessonbox}

% ================= TITLE: ROLE PLAY =================
\vspace{2em}
\begin{tcolorbox}[colback=boxframe, colframe=boxframe, rounded corners, boxrule=0pt, arc=3mm, left=2mm, top=2mm, bottom=2mm]
    \Large \bfseries \color{white} 🗣️ 7. Waktunya Role Play! (Praktik Ucapkan)
\end{tcolorbox}
\vspace{1em}

\begin{convobox}{55}
    \noindent
    \jpkana{【店員】いらっしゃいませ、}{[Tenin] Irasshaimase,} 
    \jpword{なんめい}{何名}{nan-mei}\jpword{さま}{様}{sama} \jpkana{ですか？}{desu ka?}

    \vspace{1.5em}
    \noindent{\Large \color{articolor} \textbf{Arti:} [Pelayan] Selamat datang, untuk berapa orang?}
\end{convobox}

\begin{convobox}{56}
    \noindent
    \jpkana{【あなた】}{[Anata]} \jpword{ひとり}{1人}{Hitori} \jpkana{です。}{desu.}

    \vspace{1.5em}
    \noindent{\Large \color{articolor} \textbf{Arti:} [Kamu] 1 orang. (Jangan lupa angkat 1 jari!)}
\end{convobox}

\begin{convobox}{57}
    \noindent
    \jpkana{【店員】こちらへ}{[Tenin] Kochira e} \jpkana{どうぞ。}{douzo.}

    \vspace{1.5em}
    \noindent{\Large \color{articolor} \textbf{Arti:} [Pelayan] Silakan ke sebelah sini.}
\end{convobox}

\begin{convobox}{58}
    \noindent
    \jpkana{【あなた】とんこつラーメン}{[Anata] Tonkotsu raamen} 
    \jpkana{お}{o}\jpword{ねが}{願}{nega}\jpkana{いします。}{ishimasu.}

    \vspace{1.5em}
    \noindent{\Large \color{articolor} \textbf{Arti:} [Kamu] Tolong Ramen Tonkotsu.}
\end{convobox}

\begin{convobox}{59}
    \noindent
    \jpkana{【店員】}{[Tenin] }\jpword{めん}{麺}{Men} \jpkana{の}{no} 
    \jpkana{かたさは}{katasa wa} 
    \jpkana{どうなさいますか？}{dou nasaimasu ka?}

    \vspace{1.5em}
    \noindent{\Large \color{articolor} \textbf{Arti:} [Pelayan] Tingkat kekerasan mi-nya bagaimana?}
\end{convobox}

\begin{convobox}{60}
    \noindent
    \jpkana{【あなた】}{[Anata]} \jpword{ふつう}{普通}{Futsuu} \jpkana{で}{de} 
    \jpkana{お}{o}\jpword{ねが}{願}{nega}\jpkana{いします。}{ishimasu.}

    \vspace{1.5em}
    \noindent{\Large \color{articolor} \textbf{Arti:} [Kamu] Tolong yang biasa (futsuu).}
\end{convobox}

\begin{convobox}{61}
    \noindent
    \jpkana{【店員】かしこまりました。}{[Tenin] Kashikomarimashita.} 
    \jpkana{トッピングは}{Toppingu wa} 
    \jpkana{いかがですか？}{ikaga desu ka?}

    \vspace{1.5em}
    \noindent{\Large \color{articolor} \textbf{Arti:} [Pelayan] Baik. Bagaimana dengan topping?}
\end{convobox}

\begin{convobox}{62}
    \noindent
    \jpkana{【あなた】}{[Anata]} \jpword{あじたま}{味玉}{Ajitama} \jpkana{と}{to} 
    \jpkana{ねぎ}{negi} \jpkana{トッピング}{toppingu} 
    \jpkana{お}{o}\jpword{ねが}{願}{nega}\jpkana{いします。}{ishimasu.}

    \vspace{1.5em}
    \noindent{\Large \color{articolor} \textbf{Arti:} [Kamu] Tolong topping Ajitama (telur) dan Negi (daun bawang).}
\end{convobox}

\begin{convobox}{63}
    \noindent
    \jpkana{【店員】かしこまりました。}{[Tenin] Kashikomarimashita.}

    \vspace{1.5em}
    \noindent{\Large \color{articolor} \textbf{Arti:} [Pelayan] Baik, saya mengerti pesanan Anda.}
\end{convobox}

% --- SUMMARY BOX ---
\vspace{2em}
\begin{tcolorbox}[
    colback=blue!5!white,
    colframe=blue!80!black,
    boxrule=3pt,
    arc=5mm,
    title={\huge\bfseries \sffamily 🌟 Rangkuman Misi Ramen 🌟},
    fonttitle=\color{white},
    attach boxed title to top center={yshift=-4mm},
    boxed title style={colback=blue!80!black, rounded corners},
    left=5mm, right=5mm, top=7mm, bottom=5mm
]
\Large \textbf{Omedetou! (Selamat!) Kamu siap makan ramen sepuasnya di Jepang!}
\begin{itemize}
    \item \textbf{Pesan Cepat:} Tunjuk gambar dan bilang \textit{Kore, onegaishimasu!}
    \item \textbf{Menghitung Orang:} Hitori (1), Futari (2), San-nin (3), Yo-nin (4).
    \item \textbf{Menghitung Porsi:} Hitotsu (1), Futatsu (2).
    \item \textbf{Custom Ramen:} \textit{Katame} (Keras/Kenyal), \textit{Futsuu} (Biasa), \textit{Yawarakame} (Lembek/Lunak).
    \item \textbf{Mesin Tiket (Kenbaiki):} Masukkan uang DULU, baru tekan tombol!
\end{itemize}
Jangan lupa berlatih pesan dengan orang tuamu di rumah ya! \textit{Oishii raamen o tabete kudasai ne!} (Silakan makan ramen yang enak ya!).
\end{tcolorbox}